\subsubsection{Tree DP (reroot)}
\begin{lstlisting}[style=mystyle, language=C++]
// TC: O(N)
// usa: para cada nodo, calcular alguna propiedad asumiendo que ese nodo es la raiz
#include <bits/stdc++.h>
using namespace std;

struct TreeDP{
	int n;
	vector<vector<int>> adj;
	vector<long long> dp, up, ans;
	vector<int> sz;
	TreeDP(int n): n(n){
		adj.assign(n, {});
		dp.assign(n, 0);
		up.assign(n, 0);
		ans.assign(n, 0);
		sz.assign(n, 0);
	}
	void addEdge(int u, int v){
		adj[u].push_back(v);
		adj[v].push_back(u);
	}
	// 1st pass: compute dp[v] as the answer for subtree of v (rooted at 0)
	void dfs1(int v, int p){
		sz[v] = 1;
		for(int u : adj[v]){
			if(u == p) continue;
			dfs1(u, v);
			sz[v] += sz[u];
			dp[v] += dp[u] + sz[u]; // example - sum of distances from u
		}
	}
	// 2nd pass: reroot - propagate contributions
	void dfs2(int v, int p){
		ans[v] = dp[v] + up[v];
		for(int u : adj[v]){
			if(u == p) continue;
			// when rerooting at u, u's contribution comes out, v's grows
			up[u] = up[v] + dp[v] - (dp[u] + sz[u]) + (n - sz[u]);
			dfs2(u, v);
		}
	}
	vector<long long> solve(int root = 0){
		dfs1(root, -1);
		dfs2(root, -1);
		return ans;
	}
};

int main(){
	TreeDP tdp(4);
	tdp.addEdge(0,1); tdp.addEdge(1,2); tdp.addEdge(1,3);
	auto ans = tdp.solve(0);
	for(long long x : ans) cout << x << " ";  // sum of distances from each node
	cout << "\n";
	return 0;
}
\end{lstlisting}
